% ============================================================
% 准静态电磁场（子文件，只放内容：不重复 \usepackage、不改字号）
% 对应教材：冯慈璋·马西奎《工程电磁场导论》第 5 章 §5-1 ～ §5-7
% 本章课件无单独讲义，内容以教材为准
% ============================================================

% formulaBox 由 3.矢量分析.tex 定义；这里做一次“未定义才定义”的保护。
\makeatletter
\@ifundefined{formulaBox}{%
  \newenvironment{formulaBox}{%
    \par\begin{adjustbox}{max width=\linewidth}%
    \begin{tcolorbox}[colback=blue!8!white,colframe=transparent,boxrule=0pt,halign=center,boxsep=2pt,left=4pt,right=4pt,top=3pt,bottom=3pt]%
  }{%
    \end{tcolorbox}\end{adjustbox}\par
  }%
}{}
\makeatother

\chapter{准静态电磁场}

时变场中，若感应电场远小于库仑电场（可忽略 $\partial\vec{B}/\partial t$），称\textbf{电准静态
场}（EQS）；若位移电流远小于传导电流（可忽略 $\partial\vec{D}/\partial t$），称\textbf{磁
准静态场}（MQS）。两者统称\textbf{准静态电磁场}：\textbf{属时变场，却保留静态场与源“瞬时
对应”的性质}（略去了波动性与推迟效应）。电气工程中的问题绝大部分属 MQS。

\section{电准静态场与磁准静态场}

\subsection{电准静态场（EQS）}

忽略电磁感应项 $\partial\vec{B}/\partial t$，即 $\vec{E}_j\ll\vec{E}_C$ 时，电场无旋：
$\nabla\times\vec{E}\approx\nabla\times\vec{E}_C=0$。其基本方程组为

\begin{formulaBox}
\[
\nabla\times\vec{E}=0,
\qquad
\nabla\cdot\vec{D}=\rho_f,
\qquad
\nabla\times\vec{H}=\vec{J}_C+\frac{\partial\vec{D}}{\partial t},
\qquad
\nabla\cdot\vec{B}=0
\]
\end{formulaBox}

与静态情形相比，\textbf{只有磁场方程变了}（计入位移电流产生的磁场），电场方程与静电场完全
相同；只是 $\vec{E}$、$\vec{D}$ 随时间变化，但与源\textbf{瞬时对应}。故已知电荷分布即可用静电
场公式求 $\vec{E}$、$\vec{D}$，再由全电流定律求 $\vec{H}$。典型对象：高压电场、电子器件附近
电场、低频交流平板电容器内的电场。

放宽情形：若 $\vec{E}_j$ 不小但其\textbf{旋度}很小，式仍成立，也可按 EQS 处理（如低频电感
线圈导线中的电场可按恒定电场处理）。

\subsection{磁准静态场（MQS）}

忽略位移电流项 $\partial\vec{D}/\partial t$，磁场按恒定磁场处理：

\begin{formulaBox}
\[
\nabla\times\vec{H}=\vec{J}_C,
\qquad
\nabla\cdot\vec{B}=0,
\qquad
\nabla\times\vec{E}=-\frac{\partial\vec{B}}{\partial t},
\qquad
\nabla\cdot\vec{D}=\rho_f
\]
\end{formulaBox}

与静态情形相比，\textbf{只有电场方程变了}（计入电磁感应），磁场方程与恒定磁场相同。与恒定
磁场一样，$\vec{B}=\nabla\times\vec{A}$、$\vec{E}=-\partial\vec{A}/\partial t-\nabla\varphi$，
且在库仑规范 $\nabla\cdot\vec{A}=0$ 下

\begin{formulaBox}
\[
\nabla^{2}\vec{A}=-\mu\vec{J}_C,
\qquad
\nabla^{2}\varphi=-\frac{\rho_f}{\varepsilon}
\]
\end{formulaBox}

即 $\vec{A}$、$\varphi$ 随时间变却遵循静态场（泊松方程）的规律：已知电流、电荷分布即可用静态
公式求位。即可略去场的波动性，认为场与源具有类似静态场的“瞬时”对应关系，故 MQS 又称
\textbf{似稳场}。典型对象：电机、变压器、感应加热、磁悬浮的磁场，及电工中的涡流问题。

\subsection{似稳条件}

\textbf{（1）导体中的时变场}：若 $\dfrac{\omega\varepsilon}{\gamma}\ll1$，位移电流可略，导体内
时变场按 MQS 处理。满足此条件的导体称\textbf{良导体}，其中的 MQS 又称\textbf{涡流场}；对纯
金属该条件一直到紫外波段都成立。

\textbf{（2）理想介质中的时变场}：与源相距 $R$ 处的相位因子为 $e^{-j\beta R}$，要使其近似为 1，
须
\[
\beta R=\frac{2\pi R}{\lambda}\ll1
\quad\Longleftrightarrow\quad
R\ll\lambda
\]
满足该条件的区域称\textbf{近区}（似稳区），其中场与源同步变化、无相位移。两条件统称
\textbf{MQS 近似条件}（似稳条件）。它给出重要判据：\textbf{能否用“路”的方法处理，取决于尺寸
与波长之比，而非绝对尺寸或频率高低}——工频 $50\,\mathrm{Hz}$ 时 $\lambda\approx6000\,\mathrm{km}$，
只有数百公里的输电才需计波动；$3\,\mathrm{GHz}$ 时 $\lambda=10\,\mathrm{cm}$，手掌大的系统
就须计波动。似稳区是\textbf{相对}概念。

\corepoint{\textbf{EQS 丢掉的是“变化磁场激发电场”这条通路（$\partial\vec{B}/\partial t$），
保留位移电流；MQS 丢掉的是“变化电场激发磁场”这条通路（$\partial\vec{D}/\partial t$），保留
电磁感应。}两者\textbf{不能同时}用，否则退化为静态场、丢掉全部耦合与波动信息。}

\section{磁准静态场和电路}

MQS 方程是交流电路的理论基础，基尔霍夫两条定律均可由它导出。

\textbf{（1）基尔霍夫电流定律。}对 $\nabla\times\vec{H}=\vec{J}_C$ 取散度得
$\oint_S\vec{J}_C\cdot \mathrm{d}\vec{S}=0$，即 MQS 中传导电流\textbf{连续}；用于包围节点的
闭合面即得 $\sum i_k=0$。

\textbf{（2）基尔霍夫电压定律。}由 $\vec{J}_C=\gamma(\vec{E}+\vec{E}_e)$ 得
$\vec{E}=\vec{J}_C/\gamma-\vec{E}_e$，代入
$\vec{E}=-\partial\vec{A}/\partial t-\nabla\varphi$ 沿导线由 $A$ 到 $B$ 积分，四项依次对应
电动势、电阻压降、电感压降与电容电压：

\begin{formulaBox}
\[
e(t)=\int_{A}^{B}\vec{E}_e\cdot \mathrm{d}\vec{l}
=i\left(R+r+R_i\right)+L\frac{\mathrm{d}i}{\mathrm{d}t}+u_C
\]
\end{formulaBox}
即“沿闭合回路的电动势等于各段电压降之和”。可见电路理论不过是 $R\ll\lambda$、
$\omega\varepsilon\ll\gamma$ 条件下麦克斯韦理论的近似。

\corepoint{时变场中“电压”只在\textbf{同一横截面上的两点间}才有意义：在纵剖面上取闭合路径，
$\oint_l\vec{E}\cdot \mathrm{d}\vec{l}\neq0$（所围面上有磁通）；只有取在同一横截面内的闭合
路径，才因所围面无磁通而等于零。}

\section{电准静态场与电荷弛豫}

在 $\gamma$、$\varepsilon$ 均匀的导体内，由电荷守恒定律、$\vec{J}_C=\gamma\vec{E}$ 与
$\nabla\cdot\vec{E}=\rho/\varepsilon$ 得

\begin{formulaBox}
\[
\frac{\partial\rho}{\partial t}+\frac{\gamma}{\varepsilon}\rho=0
\quad\Longrightarrow\quad
\rho(t)=\rho_0e^{-t/\tau},
\qquad
\tau=\frac{\varepsilon}{\gamma}
\]
\end{formulaBox}

自由电荷体密度按指数规律衰减，$\tau$ 称\textbf{弛豫时间}。良导体的 $\tau$ 极小，故可认为
\textbf{良导体内部无自由电荷积累}（$\rho=0$），电荷只能分布在表面——与静电场结论一致，但这里
是\textbf{动态}得到的。

EQS 下仍可定义 $\vec{E}=-\nabla\varphi$，代入高斯定律得

\begin{formulaBox}
\[
\nabla^{2}\varphi=-\frac{\rho}{\varepsilon}
=-\frac{\rho_0}{\varepsilon}e^{-t/\tau},
\qquad
\varphi(\vec{r},t)=\varphi(\vec{r},0)\,e^{-t/\tau}
\]
\end{formulaBox}
（后式为无限大均匀导电媒质中的解。）即电位分布按同一指数规律衰减，快慢同取决于 $\tau$；对点
电荷 $\varphi=\dfrac{q_0}{4\pi\varepsilon R}e^{-t/\tau}$，而磁场处处为零。

\textbf{分片均匀导体。}电荷趋于聚积在两导体分界面上，界面两侧须同时满足衔接条件与连续性条件：
\[
E_{1t}=E_{2t},
\qquad
D_{2n}-D_{1n}=\sigma,
\qquad
J_{1n}=J_{2n}\ \Longrightarrow\ \gamma_1E_{1n}=\gamma_2E_{2n}
\]
以双层有损介质平板电容器接直流电压为例：$t=0^+$ 时界面电荷来不及建立，电场按\textbf{介电
常数}分配（如两层理想介质）；$t\to\infty$ 稳态仅有传导电流，电场按\textbf{电导率}分配（如两层
导电媒质）；其间按指数过渡。工程结论：低频（工频）交流下若位移电流远大于漏电流，电场按
$\varepsilon$ 分布，属静电场问题；直流下按 $\gamma$ 分布，属恒定电场问题，两种稳态一般不同。

\corepoint{判据两种写法等价：$\omega\varepsilon/\gamma\ll1$ 与 $\omega\tau\ll1$
（$\tau=\varepsilon/\gamma$），前者从“位移电流/传导电流”比值得出，后者从“周期/弛豫时间”得出。
$\gamma$ 是与 $\omega$、$\varepsilon$ \textbf{相对}比较的：同一材料低频是良导体，到光频可能
变成良介质。}

\section{集肤效应}

MQS 下略去位移电流，对 $\nabla\times\vec{H}=\vec{J}_C$ 取旋度并利用矢量恒等式与物性方程，得

\begin{formulaBox}
\[
\nabla^{2}\vec{H}=\mu\gamma\frac{\partial\vec{H}}{\partial t},
\quad
\nabla^{2}\vec{E}=\mu\gamma\frac{\partial\vec{E}}{\partial t},
\quad
\nabla^{2}\vec{J}_C=\mu\gamma\frac{\partial\vec{J}_C}{\partial t}
\]
\[
\text{正弦稳态：}\quad
\nabla^{2}\dot{\vec{H}}=j\omega\mu\gamma\dot{\vec{H}},
\quad
\nabla^{2}\dot{\vec{E}}=j\omega\mu\gamma\dot{\vec{E}},
\quad
\nabla^{2}\dot{\vec{J}}_C=j\omega\mu\gamma\dot{\vec{J}}_C
\]
\end{formulaBox}
它们也称\textbf{涡流方程}或\textbf{磁扩散方程}（时域形式描述磁场扩散过程）。取 $x>0$ 为导体、
正弦电流沿 $y$ 方向且只随 $x$ 变，则
$\mathrm{d}^{2}\dot{J}_y/\mathrm{d}x^{2}=k^{2}\dot{J}_y$，舍去随 $x$ 增大的项，得

\begin{formulaBox}
\[
k=\frac{1+j}{\delta},
\qquad
\dot{J}_y=\dot{J}_0e^{-kx}=\dot{J}_0e^{-x/\delta}e^{-jx/\delta},
\qquad
\dot{E}_y=\dot{E}_0e^{-kx},
\qquad
\dot{H}_z=-\frac{jk}{\omega\mu}\dot{E}_0e^{-kx}
\]
\[
\delta=\frac{1}{\alpha}=\sqrt{\frac{2}{\omega\mu\gamma}}
=\frac{1}{\sqrt{\pi f\mu\gamma}}
\]
\end{formulaBox}

$\delta$ 为\textbf{透入深度}：场量振幅衰减到表面值的 $1/e$（约 36.8\%）所经距离。场量振幅沿纵深
按 $e^{-x/\delta}$ 衰减且\textbf{相位随之改变}，愈深愈小——这种场量集中在导体表面附近的现象即
\textbf{集肤效应}。由 $\dot{H}_z/\dot{E}_y$ 还可见：导体内磁场\textbf{落后}电场 $45^{\circ}$、
幅值比为 $\sqrt{\omega\mu/\gamma}$（与理想介质中“同相、比值 $\sqrt{\mu/\varepsilon}$”不同）。

要点：$f$、$\gamma$、$\mu$ 越大则 $\delta$ 越小、集肤越显著（铜在 $50\,\mathrm{Hz}$ 时
$\delta\approx9.4\,\mathrm{mm}$）；经 $4.6\delta$ 振幅降为 $1\%$；深度大于数个 $\delta$ 后场量
近似为零但\textbf{不等于零}；$\delta$ 由无限大平面导体导出，只要表面曲率半径远大于 $\delta$
仍可近似使用。工程应用：高频表面淬火；高频线圈用空心铜管或镀银线。

\section{涡流及其损耗}

大块导体处在变化磁场中时，内部感应出\textbf{在大块导体内部自成闭合回路、呈旋涡状流动}的电流，
称\textbf{涡旋电流}（涡流）。它有两个效应：\textbf{热效应}（损耗发热）与\textbf{去磁效应}
（涡流磁场反抗外加场的变化）。工业上用热效应加热、冶炼金属，用去磁效应制电磁阻尼、电磁闸；
变压器铁心中的涡流损耗则须设法减小。MQS 下涡流场同样遵守上节的涡流方程。

把铁心叠片视为厚度 $d$ 的薄导电平板，外磁场沿板面法向、场量与 $y$、$z$ 无关，取板中心为
$x=0$，则

\begin{formulaBox}
\[
\dot{H}_z(x)=\dot{H}_0\frac{\operatorname{ch}(kx)}
{\operatorname{ch}(kd/2)},
\qquad
\dot{J}_y(x)=-\dot{H}_0k\,
\frac{\operatorname{sh}(kx)}{\operatorname{ch}(kd/2)},
\qquad
k=\frac{1+j}{\delta}
\]
\end{formulaBox}

磁场在\textbf{板中心最小}（涡流去磁），涡流密度\textbf{反对称于中心}、中心为零、表面最大，
愈深愈小，同样呈集肤效应。对电工钢片（$\mu_r\sim1000$、$d=0.5\,\mathrm{mm}$），$50\,\mathrm{Hz}$
时 $B$ 可视为沿截面均匀；$2000\,\mathrm{Hz}$ 时板中 $B$ 已远小于表面，须改用更薄钢片，更高频
则用粉状材料压制的铁心。体积 $V$ 中的平均损耗及低频近似为

\begin{formulaBox}
\[
P=\int_{V}\frac{1}{\gamma}\left|\dot{J}_y\right|^{2}\mathrm{d}V,
\qquad
P=\frac{1}{12}\gamma\omega^{2}d^{2}B_{av}^{2}V
\quad\left(\frac{d}{\delta}\ll1\right)
\]
\end{formulaBox}
$B_{av}$ 为磁感应强度在板厚上的平均值。减损三途径：\textbf{减小厚度}（$P\propto d^{2}$，故用
彼此绝缘的硅钢片叠装，片间绝缘切断涡流回路）、\textbf{减小电导率}（掺硅提高电阻率）、
\textbf{降低频率或磁密}。

\corepoint{低频下 $P\propto f^{2}B_{av}^{2}d^{2}\gamma$。三处反直觉：（1）$\gamma$ 越\textbf{大}
损耗越\textbf{大}，故靠“掺硅降 $\gamma$”减损，而非提高导电性；（2）$P\propto d^{2}$ 只在
$d\ll\delta$ 时成立，高频下须用 $P=\int|J|^{2}/\gamma\,\mathrm{d}V$ 精算；（3）硅钢片若片间
不绝缘，涡流回路重新闭合，减损效果全部丧失。}

\section{导体的交流内阻抗}

交流下集肤效应使电流与场集中于表面附近，\textbf{实际载流截面减小}，故电阻与内电感不同于直流。
设通有总电流 $\dot{I}$（有效值），等效参数 $Z=R+jX$：一方面 $\dot{I}\dot{I}^{*}Z=I^{2}(R+jX)$，
另一方面由坡印亭定理流入导体的复功率为
$-\oint_S(\dot{\vec{E}}\times\dot{\vec{H}}^{*})\cdot \mathrm{d}\vec{S}$，两者相等得

\begin{formulaBox}
\[
Z=\frac{-\oint_{S}\left(\dot{\vec{E}}\times\dot{\vec{H}}^{*}\right)
\cdot \mathrm{d}\vec{S}}{I^{2}}=R+jX
\]
\[
R=\frac{1}{I^{2}}\mathrm{Re}\left[-\oint_{S}
\left(\dot{\vec{E}}\times\dot{\vec{H}}^{*}\right)\cdot \mathrm{d}\vec{S}\right]
\]
\[
X=\frac{1}{I^{2}}\mathrm{Im}\left[-\oint_{S}
\left(\dot{\vec{E}}\times\dot{\vec{H}}^{*}\right)\cdot \mathrm{d}\vec{S}\right]
\]
\end{formulaBox}
$Z$ 只计及\textbf{导体内部}场引起的阻抗，故称\textbf{交流内阻抗}：$R$ 为交流电阻（焦耳损耗），
$X=\omega L_i$ 为内电抗（体内磁场储能）。

\textbf{（1）半无限大导体。}取底面积 $h\times a$ 的柱体，由 $\dot{I}=a\gamma\dot{E}_0/k$ 与坡印亭
矢量只在 $x=0$ 面有通量，得

\begin{formulaBox}
\[
Z=R+jX=\frac{h}{a\gamma\delta}(1+j),
\qquad
R=\frac{h}{a\gamma\delta},
\qquad
L_i=\frac{X}{\omega}=\frac{h}{a\gamma\delta\omega}
\]
\end{formulaBox}
意义：有效值 $I$ 的交流电流给出的有功功率，等于同值直流电流\textbf{集中在厚度 $\delta$ 的层内}
的功率——透入深度的另一物理含义；因 $\delta$ 很小，$Z$ 也称\textbf{表面阻抗}。

\textbf{（2）圆导线。}半径 $a\gg\delta$ 时可视为厚度无限大、宽为周长 $2\pi a$ 的平面导体：

\begin{formulaBox}
\[
R=\frac{1}{2\pi a\gamma\delta},
\qquad
R_{dc}=\frac{1}{\gamma\pi a^{2}},
\qquad
\frac{R}{R_{dc}}=\frac{a}{2\delta}
\]
\end{formulaBox}
即交流电阻相当于厚度为 $\delta$ 的圆管的直流电阻。故同截面下减小交流电阻的唯一办法是\textbf{增大
表面面积}——用相互绝缘的多股细线（绞线），高频时用空心铜管。

\corepoint{三处易错：（1）$\dfrac{R}{R_{dc}}=\dfrac{a}{2\delta}$ 要求 $a\gg\delta$（强集肤）；
低频时 $\delta\gg a$，有 $R\approx R_{dc}$，不可套用。（2）$L_i=X/\omega$ 随频率变化（$f$ 升高
$\delta$ 减小、$L_i$ 减小），非常数。（3）式中 $I$ 是\textbf{有效值}。}

\section{邻近效应和磁屏蔽}

\textbf{（1）邻近效应。}相互靠近的导体通交变电流时，每一导体还处于其它导体电流的场中，电流
分布因而改变，称\textbf{邻近效应}；\textbf{频率愈高、靠得愈近愈显著}，与集肤效应共存。二线
传输线中电流集中于两导线\textbf{相对的内侧}：能量主要经两线之间的空间以电磁波形式传给负载，
内侧能量密度大、传入导线的功率大；通以反向电流的平行汇流排，内侧面电流密度也最大。

\textbf{（2）电磁屏蔽。}用来减弱某些源在空间某区域（不含这些源）产生的电磁场的结构，称
\textbf{电磁屏蔽}，是抑制邻近效应的常用措施，屏蔽体多用铜、铝、钢。它利用\textbf{导体内涡流的
去磁作用}抵制外加场，故又称\textbf{涡流屏蔽}；有效屏蔽要求屏蔽层厚度接近透入深度的 3～6 倍：

\begin{formulaBox}
\[
h\ \geq\ (3\sim6)\,\delta,
\qquad
S=\frac{E_0}{E}\ \ \text{或}\ \ S=\frac{H_0}{H}
\]
\end{formulaBox}
$S$ 称\textbf{屏蔽系数}（无屏蔽体时该区场强与有屏蔽体时之比）。$1\,\mathrm{MHz}$ 时铝的
$\delta\approx82\,\mu\mathrm{m}$，薄铝罩即可屏蔽射频场；\textbf{高频不宜用铁磁材料}（损耗大、
发热重），而低频恒定磁场恰要靠高 $\mu$ 铁磁材料“短路”磁通来屏蔽。

\textbf{（3）磁扩散过程。}薄壁导体圆管（内半径 $a$、壁厚 $\Delta$）外突然建立轴向场 $H_0$ 时，
感应电流的去磁作用使管内 $H_i$ 不能突变，随后感应电流衰减、磁场逐渐透入，最后均匀分布：

\begin{formulaBox}
\[
\frac{\mathrm{d}H_i}{\mathrm{d}t}+\frac{1}{\tau_m}H_i=\frac{H_0}{\tau_m},
\qquad
H_i(t)=H_0\left(1-e^{-t/\tau_m}\right),
\qquad
\tau_m=\frac{\mu_0\gamma\Delta a}{2}
\]
\[
\text{正弦稳态：}\quad
\dot{H}_i=\frac{H_0}{1+j\omega\tau_m},
\qquad
\omega\gg\frac{1}{\tau_m}=\frac{2}{\mu_0\gamma\Delta a}
\ \text{时}\ H_i\ll H_0
\]
\end{formulaBox}
$\tau_m$ 称\textbf{磁扩散时间}。可见薄壁导体圆管中感应电流的作用是\textbf{去磁}的；若管内再套
一管，则内管磁场极弱——这正是多层屏蔽更有效的原因。

\corepoint{电磁屏蔽三处易错：（1）$h\gtrsim(3\sim6)\delta$ 而 $\delta\propto1/\sqrt{f}$，故
\textbf{频率越高屏蔽层越薄}；对低频（尤其直流）磁场，薄金属罩几乎无效，必须用高 $\mu$ 铁磁
材料。（2）屏蔽靠涡流去磁，屏蔽体必须\textbf{闭合、导电良好}，开缝开孔会显著降低效能。
（3）须避开\textbf{谐振}：被屏蔽场频率接近屏蔽体固有频率时效能急剧下降，甚至使防护区场强
反而增大。}
